\documentclass[10pt,a4paper]{amsart}
\usepackage[margin=19mm]{geometry}
\usepackage{amsmath,amssymb,mathtools}
\usepackage[scr=rsfs,cal=euler]{mathalfa}
\usepackage{xfrac}
\usepackage[unicode,hidelinks,bookmarks=false]{hyperref}
\newcommand{\ie}{i.e.\ }
\newcommand{\cf}{cf.\ }
\newcommand{\resp}{respectively\ }
\newcommand{\Subsection}{\subsection}
\providecommand{\nobreakdash}{\nobreak\hbox{-}}
\setlength{\emergencystretch}{2em}
\multlinegap0pt
\usepackage{amssymb}
\usepackage{algorithm}\usepackage{algpseudocode}
\newtheorem{lemma}{Lemma}
\title{Structural Parameters for Dense Temporal Graphs}
\author{Jessica Enright, Samuel D. Hand, Laura Larios-Jones, Kitty Meeks}
\date{Source: arXiv:2404.19453v2 (CC BY 4.0)}
\begin{document}
\maketitle

\noindent\emph{Source and licence.} Structural Parameters for Dense Temporal Graphs. Version arXiv:2404.19453v2 (CC BY 4.0), distributed by arXiv under the Creative Commons Attribution 4.0 International licence, http://creativecommons.org/licenses/by/4.0/, as recorded in the arXiv licence metadata for this exact version and shown on the abstract page https://arxiv.org/abs/2404.19453v2. Full licence notice: \texttt{licenses/arxiv-2404.19453-CC-BY-4.0.txt}.\par\smallskip

\noindent\textit{Editorial scope.} Counted unit: the article's lemma lem:arbord1 (source0.tex lines 768--780). The article's complete original proof of it is delivered with it (source0.tex lines 781--799, one \texttt{proof} environment of 19 lines). No mathematics is written, repaired or re-derived: the statement and the proof are the article's own lines, in the article's own notation, and no retained line is substituted or re-flowed.  Retained uncounted as the source-local context the proof needs: lines 681--684; lines 697--710.  Nothing was omitted from any retained span, so the package carries no omission record.  No retained line contains a citation, so the wrapper carries no bibliography and no bibliography entry.  No retained line contains a figure environment, an image-inclusion command or drawing code, so no image is delivered and the package declares no asset dependency.  Editorial formatting only, in addition: an AMS \texttt{amsart} preamble that restates the article's own declarations and macros used by the retained lines, namely the environments \texttt{lemma} on separate counters, together with the standard packages those lines need.  The retained environments print in this document as Lemma 1 in order of appearance, whereas the article prints the same items with the numbering of its own full document.  The complete native source of the article as distributed in the arXiv e-print for this exact version is bundled unchanged as \texttt{sources/158385/source0.tex}.
\begin{lemma}
\label{lem:settimelin1}
Let $S$ be a successful strategy for $(G, \lambda)$. Suppose that there is some timestep $t_1 < |S|$ and strategy $R$ with $|R|=|S|$ such that $B_{t_1}(S) \subseteq B_{t_1}(R)$ and in every timestep after $t_1$, $R$ places a fire at the same vertex as $S$. Then $R$ is also successful.
\end{lemma}

\begin{lemma}
    \label{lem:alreadyburn1}
    Let $(G, \lambda)$ be a temporal graph with temporal neighbourhood partition $(X_i)_{i\in I}$. Now let $S$ be a strategy that burns this graph, and let $u$ be a vertex at which $S$ places a fire in a timestep $t_1$, and $X_i$ be the class to which this vertex belongs. Fix a timestep $t_2 > t_1$, and let $S'$ be a strategy which plays as follows until timestep $t_2$:

    \[
    s'_t=\begin{cases}
        s_t&\text{if }t<t_1\\
        s_{t+1}&\text{if }t_1\leq t < t_2\\
        u&\text{if }t=t_2
    \end{cases}
    \]

    If there exists a vertex $w \in X_i$ which is burning before the end of timestep $t_1$ when $S'$ is played, we have that $B_{t_2}(S) \subseteq B_{t_2}(S')$.
\end{lemma}

\begin{lemma}
    \label{lem:arbord1}
    Let $(G, \lambda)$ be a temporal graph, and $S$ a successful strategy such that the first $|C(S)|$ fires placed by $S$ are placed in distinct classes from the temporal neighbourhood partition. Let $f : [|C(S)|+1,|S|] \to [|C(S)|+1,|S|]$ be any bijection.

   Then the strategy $S'$ given by

    \[s'_t=\begin{cases}
        s_t&\text{if }t\leq|C(S)|\\
        s_{f(t)}&\text{otherwise}
    \end{cases}\]

    is successful.
\end{lemma}

\begin{proof}
    We begin by showing that $S'$ is in fact successful. We argue this by contradiction. Suppose there is a vertex $v$ which is not burning once all the moves in $S'$ are played. Then $v$ must be a vertex on which no fire is placed in either $S$ or $S'$. This implies that the fire spreads to $v$ from another vertex; denote by $w$ the vertex from which the fire spreads. Assume without loss of generality that $w$ is the first in its class to be burned. Then $w$ must be burned in the first $C(S)$ timesteps, and so the fire spreads to $v$ in both strategies. This contradicts the assumption that $v$ is never burned, showing that $S'$ is a successful strategy.

    % Let $(G, \lambda)$, along with a strategy $S$ and bijection $f$ be a counterexample.

    % Then let $R$ be a successful strategy with $|R| \leq |S|$, $C(R) = C(S)$, and $r_t = s_t$ for all $t \leq |C(S)|$. Furthermore, assume that $R$ is the strategy minimal in the timestep $t_2$ such that, for every timestep $t \geq t_2$, $r_t=s_{f(t)}$. (Note that it is possible that $t_2=|R\,|+1$, and there is no terminal sub-sequence on which $R$ agrees with the permutation of $S$.) Let $t_1$ be the timestep in which $R$ places a fire at $s_{f(t_2-1)}=s'_{t_2-1}$.

    % Now let $R'$ be the strategy that makes moves as follows:
    % \[
    % r'_t = \begin{cases}
    %     r_t&\text{if }t<t_1\text{ or }t\geq t_2\\
    %     r_{t+1}&\text{if }t_1\leq t < t_2-1\\
    %     s_{f(t_2-1)}&\text{if }t = t_2-1.
    % \end{cases}
    % \]

    % Then, as $R'$ places a fire at a vertex in every class in $C(S)$ prior to timestep $t_1$, by Lemma~\ref{lem:alreadyburn1} we have that $B_{t_2-1}(R) \subseteq B_{t_2-1}(R')$. Then, by Lemma~\ref{lem:settimelin1}, $R'$ burns $(G, \lambda)$ in the same or less time as $R$. This contradicts the assumption that $R$ was minimal in the timestep $t_2$, so no such counterexample $(G, \lambda)$, strategy $S$, and bijection $f$ can exist. 

\end{proof}

\end{document}
